\documentclass{article}
\usepackage{amsmath,amssymb,amsthm}
\usepackage{algorithm}
\usepackage[noend]{algpseudocode}
\usepackage{bm}
\usepackage{hyperref}
\usepackage{enumitem}
\usepackage{color}
\newcommand{\E}{\mathbb{E}}
\newcommand{\R}{\mathbb{R}}
\newcommand{\norm}[1]{\left\|#1\right\|}
\newcommand{\ip}[2]{\langle #1,#2\rangle}
\newcommand{\step}[1]{\textbf{[Step #1]}}
\begin{document}

\section*{Differentially Private Stochastic Gradient Descent (DP‑SGD) with Gaussian Noise}

\section*{Mathematical Formulation}
Let $f:\R^{d}\to\R$ be the empirical risk
\[
f(w)=\frac{1}{n}\sum_{i=1}^{n} \ell(w;z_i),
\]
where each $\ell(\cdot;z_i)$ is convex and $L$–smooth.  
At iteration $t\in\{0,1,\dots,T-1\}$ we perform

\[
\boxed{
\begin{aligned}
g_t &:= \nabla \ell(w_t;z_{i_t}) + \xi_t,\\
w_{t+1} &:= w_t - \eta_t g_t ,
\end{aligned}}
\tag{1}
\]

where  

\begin{itemize}
\item $i_t$ is drawn uniformly from $\{1,\dots,n\}$ (mini‑batch size $=1$ for simplicity);
\item $\xi_t\sim\mathcal{N}\!\big(0,\sigma^{2}I_{d}\big)$ is independent Gaussian noise added for \emph{(ε,δ)}‑ differential privacy;
\item $\eta_t>0$ is the learning‑rate schedule (e.g. $\eta_t=\eta/\sqrt{t+1}$).
\end{itemize}

The algorithm tracks the privacy budget via the moments accountant; this bookkeeping is omitted from the mathematical derivation but assumed to be respected by the choice of $\sigma$ and $T$.

\section*{Pseudocode}
\begin{algorithm}[H]
\caption{DP‑SGD with Gaussian Noise}
\begin{algorithmic}[1]
\State \textbf{Input:} data $\{z_i\}_{i=1}^n$, initial point $w_0\in\R^{d}$, learning‑rate schedule $\{\eta_t\}$, noise scale $\sigma>0$, number of iterations $T$.
\For{$t=0$ \textbf{to} $T-1$}
    \State Sample an index $i_t\sim\text{Uniform}\{1,\dots,n\}$.
    \State Compute the (clipped) stochastic gradient $ \tilde g_t:=\nabla\ell(w_t;z_{i_t})$.
    \State Sample noise $\xi_t\sim\mathcal N(0,\sigma^{2}I_d)$.
    \State Form the private gradient $g_t := \tilde g_t + \xi_t$.
    \State Update $w_{t+1} \gets w_t - \eta_t g_t$.
\EndFor
\State \textbf{Output:} $w_T$ (or the averaged iterate $\bar w_T:=\frac1T\sum_{t=0}^{T-1}w_t$).
\end{algorithmic}
\end{algorithm}

\section*{Dry Run Example}
Consider the one‑dimensional quadratic $f(w)=(w-3)^2$, thus $\nabla f(w)=2(w-3)$.  
Take $w_0=0$, constant step size $\eta=0.1$, noise variance $\sigma^{2}=0.01$ and use the following noise realisations (drawn once for the illustration):
\[
\xi_0=0.05,\quad \xi_1=-0.02,\quad \xi_2=0.03 .
\]

\begin{center}
\begin{tabular}{c|c|c|c|c}
$t$ & $w_t$ & $\nabla f(w_t)$ & $g_t=\nabla f(w_t)+\xi_t$ & $w_{t+1}=w_t-\eta g_t$\\\hline
0 & $0.000$ & $2(0-3)=-6.00$ & $-6.00+0.05=-5.95$ & $0-0.1(-5.95)=0.595$\\
1 & $0.595$ & $2(0.595-3)=-4.81$ & $-4.81-0.02=-4.83$ & $0.595-0.1(-4.83)=1.078$\\
2 & $1.078$ & $2(1.078-3)=-3.844$ & $-3.844+0.03=-3.814$ & $1.078-0.1(-3.814)=1.460$\\
\end{tabular}
\end{center}

Corresponding objective values:
\[
f(w_0)=9,\; f(w_1)= (0.595-3)^2\approx5.78,\;
f(w_2)= (1.078-3)^2\approx3.70,\;
f(w_3)= (1.460-3)^2\approx2.38 .
\]
The iterates move toward the optimum $w^{\star}=3$, albeit perturbed by the injected Gaussian noise.

\newpage
\section*{Assumptions}
\begin{enumerate}[label=\textbf{A\arabic*:},leftmargin=*]
\item \textbf{Convexity.} $f$ is convex:
\[
f(y)\ge f(x)+\ip{\nabla f(x)}{y-x},\qquad\forall x,y\in\R^{d}.
\tag{A1}
\]
\item \textbf{L‑smoothness.} There exists $L>0$ such that
\[
\norm{\nabla f(x)-\nabla f(y)}\le L\norm{x-y},\qquad\forall x,y.
\tag{A2}
\]
Equivalently,
\[
f(y)\le f(x)+\ip{\nabla f(x)}{y-x}+\frac{L}{2}\norm{y-x}^{2}.
\tag{A2'}
\]
\item \textbf{Bounded stochastic gradients.} For every sampled index $i$,
\[
\norm{\nabla\ell(w;z_i)}\le G,\qquad\forall w\in\R^{d}.
\tag{A3}
\]
\item \textbf{Gaussian noise.} $\{\xi_t\}$ are i.i.d. $\mathcal N(0,\sigma^{2}I_d)$, independent of the data sampling. Hence
\[
\E[\xi_t]=0,\qquad \E[\norm{\xi_t}^{2}]=d\sigma^{2}.
\tag{A4}
\]
\item \textbf{Step‑size.} The learning‑rate schedule satisfies
\[
0<\eta_t\le\frac{1}{L},\qquad\text{and}\qquad \sum_{t=0}^{T-1}\eta_t =\Theta(\sqrt{T})\;\text{(e.g. }\eta_t=\eta/\sqrt{t+1}\text{)}.
\tag{A5}
\]
\end{enumerate}

\section*{Main Convergence Theorem}
\begin{theorem}[Expected Suboptimality of DP‑SGD]\label{thm:conv}
Under assumptions \textbf{A1}–\textbf{A5}, let $w^{\star}\in\arg\min_{w}f(w)$ and define the averaged iterate
\[
\bar w_T:=\frac{1}{T}\sum_{t=0}^{T-1}w_t .
\]
For a constant step size $\eta_t\equiv\eta\le\frac{1}{L}$ we have
\[
\boxed{
\E\!\big[f(\bar w_T)-f(w^{\star})\big]\le
\frac{\norm{w_0-w^{\star}}^{2}}{2\eta T}
+\frac{\eta}{2}\big(G^{2}+d\sigma^{2}\big).
}
\tag{2}
\]
Choosing $\eta=\frac{\norm{w_0-w^{\star}}}{\sqrt{T\,(G^{2}+d\sigma^{2})}}$ yields the rate
\[
\E\!\big[f(\bar w_T)-f(w^{\star})\big]=O\!\left(\frac{1}{\sqrt{T}}\right).
\tag{3}
\]
\end{theorem}

\section*{Theoretical Properties}
\begin{itemize}
\item \textbf{Stability.} The bound (2) shows that the algorithm is stable for any $\eta\le 1/L$; the additive term $\frac{\eta d\sigma^{2}}{2}$ quantifies the price paid for privacy.
\item \textbf{Hyper‑parameter Sensitivity.} The optimal $\eta$ balances the deterministic term $\frac{\norm{w_0-w^{\star}}^{2}}{2\eta T}$ against the stochastic term $\frac{\eta}{2}(G^{2}+d\sigma^{2})$. Over‑large $\eta$ inflates the variance term; too small $\eta$ slows deterministic convergence.
\item \textbf{Comparison to non‑private SGD.} Setting $\sigma=0$ recovers the classical SGD bound $\mathcal{O}(1/\sqrt{T})$ with variance $G^{2}$. The privacy penalty is linear in $\sigma^{2}$.
\item \textbf{Effect of clipping.} If gradients are clipped to norm $C$, then $G$ can be replaced by $C$, further limiting the variance contribution.
\end{itemize}

\section*{Discussion}
The theorem demonstrates that DP‑SGD retains the optimal $\mathcal{O}(1/\sqrt{T})$ convergence rate for convex, smooth objectives, up to an additive term proportional to the noise variance required for differential privacy. The result hinges only on first‑order smoothness and boundedness, matching the classical analysis of SGD while making the privacy cost explicit.

\newpage
\section*{Preliminaries (Definitions and Lemmas)}
\begin{definition}[Projection-free update]
For any vectors $a,b\in\R^{d}$ and scalar $\eta>0$,
\[
\norm{a-\eta b}^{2}= \norm{a}^{2}-2\eta\ip{a}{b}+\eta^{2}\norm{b}^{2}.
\tag{L1}
\]
\end{definition}
\begin{lemma}[Unbiasedness of noisy gradient]\label{lem:unbiased}
Under \textbf{A3} and \textbf{A4},
\[
\E[g_t\mid w_t]=\nabla f(w_t),\qquad
\E[\norm{g_t}^{2}\mid w_t]\le G^{2}+d\sigma^{2}.
\tag{L2}
\]
\end{lemma}
\begin{proof}
Since $g_t=\nabla\ell(w_t;z_{i_t})+\xi_t$ and $\E[\xi_t]=0$, the conditional expectation equals $\E[\nabla\ell(w_t;z_{i_t})]=\nabla f(w_t)$. Moreover,
\[
\E[\norm{g_t}^{2}\mid w_t]=\E\big[\norm{\nabla\ell(w_t;z_{i_t})}^{2}\mid w_t\big]
+\E[\norm{\xi_t}^{2}]
\le G^{2}+d\sigma^{2},
\]
using \textbf{A3} and \textbf{A4}. \qedhere
\end{proof}

\newpage
\section*{Main Proof (Step‑by‑Step Derivation)}
We prove Theorem~\ref{thm:conv} for a constant step size $\eta_t\equiv\eta\le 1/L$.

\noindent\textbf{Step 1.} Start from the squared distance recursion using Lemma (L1):
\[
\begin{aligned}
\norm{w_{t+1}-w^{\star}}^{2}
&= \norm{w_t-\eta g_t-w^{\star}}^{2}\\
&= \norm{w_t-w^{\star}}^{2}
-2\eta\ip{g_t}{w_t-w^{\star}}
+\eta^{2}\norm{g_t}^{2}.
\end{aligned}
\tag{4}
\]
\step{1}

\noindent\textbf{Step 2.} Rearrange (4) to isolate the inner product term:
\[
2\eta\ip{g_t}{w_t-w^{\star}}
= \norm{w_t-w^{\star}}^{2}
-\norm{w_{t+1}-w^{\star}}^{2}
+\eta^{2}\norm{g_t}^{2}.
\tag{5}
\]
\step{2}

\noindent\textbf{Step 3.} Take expectation conditioned on $w_t$ and use Lemma~\ref{lem:unbiased}:
\[
\begin{aligned}
2\eta\,\E\!\big[\ip{\nabla f(w_t)}{w_t-w^{\star}}\mid w_t\big]
&= \norm{w_t-w^{\star}}^{2}
-\E\!\big[\norm{w_{t+1}-w^{\star}}^{2}\mid w_t\big]
+\eta^{2}\E\!\big[\norm{g_t}^{2}\mid w_t\big] \\
&\le \norm{w_t-w^{\star}}^{2}
-\E\!\big[\norm{w_{t+1}-w^{\star}}^{2}\mid w_t\big]
+\eta^{2}\big(G^{2}+d\sigma^{2}\big).
\end{aligned}
\tag{6}
\]
\step{3}

\noindent\textbf{Step 4.} Apply convexity \textbf{A1} to replace the inner product:
\[
f(w_t)-f(w^{\star})\le\ip{\nabla f(w_t)}{w_t-w^{\star}}.
\tag{7}
\]
\step{4}

\noindent\textbf{Step 5.} Combine (6) and (7) and divide by $2\eta$:
\[
\begin{aligned}
f(w_t)-f(w^{\star})
&\le \frac{1}{2\eta}\Big(
\norm{w_t-w^{\star}}^{2}
-\E\!\big[\norm{w_{t+1}-w^{\star}}^{2}\mid w_t\big]\Big)
+\frac{\eta}{2}\big(G^{2}+d\sigma^{2}\big).
\end{aligned}
\tag{8}
\]
\step{5}

\noindent\textbf{Step 6.} Take full expectation over all randomness up to iteration $t$ (law of total expectation). Denote $A_t:=\E[\norm{w_t-w^{\star}}^{2}]$. Then (8) becomes
\[
\E\!\big[f(w_t)-f(w^{\star})\big]
\le \frac{A_t-A_{t+1}}{2\eta}
+\frac{\eta}{2}\big(G^{2}+d\sigma^{2}\big).
\tag{9}
\]
\step{6}

\noindent\textbf{Step 7.} Sum (9) over $t=0,\dots,T-1$:
\[
\sum_{t=0}^{T-1}\E\!\big[f(w_t)-f(w^{\star})\big]
\le \frac{A_0-A_T}{2\eta}
+ \frac{T\eta}{2}\big(G^{2}+d\sigma^{2}\big).
\tag{10}
\]
\step{7}

\noindent\textbf{Step 8.} Since $A_T\ge0$, drop it to obtain an upper bound:
\[
\sum_{t=0}^{T-1}\E\!\big[f(w_t)-f(w^{\star})\big]
\le \frac{\norm{w_0-w^{\star}}^{2}}{2\eta}
+ \frac{T\eta}{2}\big(G^{2}+d\sigma^{2}\big).
\tag{11}
\]
\step{8}

\noindent\textbf{Step 9.} Divide (11) by $T$ and use convexity of $f$ together with Jensen’s inequality:
\[
\E\!\big[f(\bar w_T)-f(w^{\star})\big]
\le \frac{1}{T}\sum_{t=0}^{T-1}\E\!\big[f(w_t)-f(w^{\star})\big]
\le \frac{\norm{w_0-w^{\star}}^{2}}{2\eta T}
+\frac{\eta}{2}\big(G^{2}+d\sigma^{2}\big).
\tag{12}
\]
\step{9}

Equation (12) is exactly the bound (2) claimed in Theorem~\ref{thm:conv}. ∎

\section*{Rate Derivation (With Constants)}
From (12) we minimise the right‑hand side with respect to $\eta$ under the constraint $\eta\le 1/L$. Ignoring the $1/L$ ceiling (the optimal $\eta$ will be smaller for large $T$), set the derivative to zero:
\[
-\frac{\norm{w_0-w^{\star}}^{2}}{2\eta^{2}T}
+\frac{1}{2}\big(G^{2}+d\sigma^{2}\big)=0
\;\Longrightarrow\;
\eta^{\star}= \frac{\norm{w_0-w^{\star}}}{\sqrt{T\,(G^{2}+d\sigma^{2})}}.
\tag{13}
\]
Plugging $\eta^{\star}$ into (12) yields
\[
\E\!\big[f(\bar w_T)-f(w^{\star})\big]
\le
\frac{\norm{w_0-w^{\star}}\sqrt{G^{2}+d\sigma^{2}}}{\sqrt{T}}.
\tag{14}
\]
Thus the expected suboptimality decays as $\mathcal{O}\!\big(1/\sqrt{T}\big)$, with a constant that grows linearly with the noise scale $\sigma$ (privacy cost) and with the gradient bound $G$.

\section*{Special Cases}
\begin{enumerate}[label=\textbf{Case \arabic*:},leftmargin=*]
\item \textbf{No privacy noise ($\sigma=0$).} Bound (14) reduces to the classic SGD rate $\mathcal{O}\big(G/\sqrt{T}\big)$.
\item \textbf{Strongly convex $f$ (parameter $\mu>0$).} Using strong convexity in place of (7) yields a linear convergence term $\big(1-\eta\mu\big)^{T}$ plus a steady‑state error $\frac{\eta (G^{2}+d\sigma^{2})}{2\mu}$.
\item \textbf{Mini‑batch of size $B$.} The variance term scales as $\frac{d\sigma^{2}}{B}$, improving the privacy‑accuracy trade‑off.
\item \textbf{Clipping at $C$.} If gradients are clipped to $\norm{\nabla\ell}\le C$, then $G$ can be replaced by $C$, providing a tighter bound when $C<G$.
\end{enumerate}

\end{document}